\documentclass[11pt,a4paper]{article}
\usepackage[utf8]{inputenc}
\usepackage[T1]{fontenc}
\usepackage{amsmath,amssymb,booktabs,geometry,graphicx}
\geometry{margin=2.5cm}
\usepackage[colorlinks=true,linkcolor=blue,citecolor=blue]{hyperref}

\title{\textbf{The Dark Dimension Framework --- Cosmic Shadow Geometry\\[4pt]
VIII: Vacuum Winding and the Origin of the Matter--Antimatter Asymmetry}}
\author{Hicham Boufourou\\ \small Independent Researcher, Brussels, Belgium}
\date{August 2026 --- v1.0}

\newcommand{\meV}{\,\mathrm{meV}}
\newcommand{\GeV}{\,\mathrm{GeV}}
\newcommand{\eV}{\,\mathrm{eV}}

\begin{document}
\maketitle

\begin{abstract}
\noindent
The condensate of Article~III carries a phase; a phase on an interval can wind.
This paper asks whether the winding of the vacuum along the fifth dimension can
bias the early universe toward matter, and answers with a computation rather
than a scenario. The mechanism belongs to the established class of spontaneous
baryogenesis (Cohen--Kaplan); what is new is the geometric origin of the
chemical potential, $\mu_w=\kappa_w\,\partial_y\Phi$, whose gradient is the
linear well of Article~II and is therefore already fixed by the series. Three
results follow. \emph{The nucleon lock:} a $\Delta B=1$ transfer operator
without a mass threshold gives the proton a nanosecond lifetime; the transfer
must pass through a state heavier than the proton, and the $B=1,\ L=1$ states
of Article~VII are exactly the required object --- the floor $m_X>m_p$ of that
article's benchmark window is hereby derived. \emph{The closure:} demanding
$\eta_B=6.1\times10^{-10}$ with the corpus numbers returns an order-one
transfer coefficient, $b_w\simeq1$--$5$, \emph{if and only if} the winding
normalisation is the meV scale of the sector itself; at the TeV wall scale the
mechanism fails by fifteen orders. The paper is therefore conditional on one
microscopic question --- who carries the winding --- and says so.
\emph{The residual:} if a trace of the bias survives today, its judges are
catalogued in an SME dictionary with inserted bounds; a narrowing theorem shows
that the only surviving route for a present-day imprint is the four-fermion $X$
operator, making the running LHCb searches of Article~VII a test of the
baryogenesis operator itself.
\end{abstract}

% ======================================================================
\section{Position of this paper}
\label{sec:position}

Earlier drafts carried the label ``DDF~0'', as if this text founded the
series. It does not: every quantitative input below --- the well slope of
Article~II, the condensate of Article~III, the shadow scale of Article~IV, the
charge classes of Article~VII --- is inherited. The honest label is
\textbf{VIII}, and the paper stands or falls with its parents.

The mechanism class is not new either, and we say so first. A chemical
potential coupled to a current, biasing equilibrium and freezing an asymmetry
at decoupling, is the \emph{spontaneous baryogenesis} of Cohen and
Kaplan~\cite{CohenKaplan}, extended by Cohen, Kaplan and
Nelson~\cite{CKN}. That class is the parent of everything below; the
contribution of this paper is confined to three items: the geometric origin of
$\mu_w$ (\S\ref{sec:bias}), the nucleon lock and its consequence for
Article~VII (\S\ref{sec:lock}), and the two-branch closure
(\S\ref{sec:closure}). [Anteriority declared.]

% ======================================================================
\section{The bias: a chemical potential from geometry}
\label{sec:bias}

\subsection{The winding charge}
The condensate of Article~III is a complex field; its phase $\theta$ lives on
the interval of Article~I. A configuration $\theta(y)$ interpolating by $2\pi
n$ across the interval carries the integer winding $n$, with current
$j_w^\mu\propto\epsilon^{\mu\nu}\partial_\nu\theta$ in the effective
two-dimensional $(t,y)$ reduction. Winding is topological: it cannot relax
locally, only through the boundaries. [Derived, standard.]

\subsection{The chemical potential}
A background gradient along $y$ acts on any current it couples to as a
chemical potential. We write
\begin{equation}
\mu_w \;=\; \kappa_w\,\partial_y\Phi\Big|_{\rm brane},
\label{eq:muw}
\end{equation}
with $\Phi$ the dark scalar whose profile Article~II derived to be
\emph{exactly linear}: slope
\begin{equation}
F \;=\; \frac{\varepsilon}{\ell} \;=\;
\frac{24\meV}{5.81\,\mu\mathrm{m}}
\;=\; 24\meV\times34.0\meV \;=\; 8.15\times10^{-4}\eV^2 .
\label{eq:slope}
\end{equation}
The gradient is therefore \emph{not a free function}: it is the capacitor of
Article~II, already fixed by the tower. What remains free is the normalisation
$\kappa_w$, of dimension (energy)$^{-1}$. Two scales are natural:
\begin{equation}
\kappa_w \sim \frac{1}{\mu}\ \ (\text{the meV scale of the sector itself})
\qquad\text{or}\qquad
\kappa_w \sim \frac{1}{M_S}\ \ (\text{the TeV wall scale}),
\label{eq:branches}
\end{equation}
giving respectively
\begin{equation}
\mu_w \;=\; \frac{F}{\varepsilon} \;=\; \frac{1}{\ell} \;=\; 34\meV
\qquad\text{or}\qquad
\mu_w \;=\; \frac{F}{M_S} \;=\; 1.4\times10^{-16}\eV .
\label{eq:muw_values}
\end{equation}
If the winding field is the condensate phase itself --- the sector's own
Goldstone --- the first normalisation is the natural one. The paper does not
choose; \S\ref{sec:closure} shows the choice is the entire verdict, and
\S\ref{sec:die} files it as the gate. [Structure derived; scale open.]

\paragraph{Assumption (declared): frozen geometry.}
The slope entering the closure is today's value~(\ref{eq:slope}). This assumes
the wall and well are at (or near) their final form at the freeze-out
temperature of \S\ref{sec:lock}, $T_f\sim0.05$--$0.2\GeV$. The assumption is
consistent --- the wall forms at $T\sim M_S\sim$~TeV, four orders above $T_f$
--- but it is an assumption, not a theorem, and a first-order history of
$\partial_y\Phi(T)$ would replace it. [Assumption.]

% ======================================================================
\section{The transfer, and the nucleon lock}
\label{sec:lock}

The bias~(\ref{eq:muw}) acts on baryons only through an operator converting
winding charge into baryon number, $\Delta w=\mp1\leftrightarrow\Delta B=\pm1$.
Sakharov's conditions are then met: $B$~violation (the transfer operator),
$C/CP$~violation (the winding direction), departure from equilibrium (the
freeze-out below).

\paragraph{Proposition L (the nucleon lock).}
\textit{The transfer operator must pass through a state heavier than the
proton.}
A contact transfer suppressed by the wall scale, $\mathcal{O}\sim
(j_w\!\cdot\!j_B)/M_S^2$, with no mass threshold, opens $p\to(\text{light
winding quantum})+\ell$ at the rate $\Gamma_p\sim m_p^5/M_S^4$, i.e.\ a proton
lifetime
\begin{equation}
\tau_p \;\sim\; 1.2\times10^{-9}\ \mathrm{s},
\end{equation}
against the observed $\tau_p>10^{34}$~yr --- an excess of fifty-one orders. The
only escape is kinematic: the minimal winding-carrying, baryon-number-carrying
excitation must satisfy $m>m_p$, closing the channel at zero temperature while
the hot plasma, at $T\gtrsim m$, crosses the threshold freely.
[Derived.]

\paragraph{Consequence for Article~VII.}
The states with exactly this charge profile --- carrying $B$ together with the
dark quantum number, heavier than the proton --- are the $X$ states of the
empty classes $(0,1,1)$ and $(1,1,1)$ of Article~VII, whose benchmark window is
$4$--$6\GeV$. Two promotions follow. First, the $X$ states, until now merely
\emph{permitted} by the Fano classification, become \emph{necessary} to this
mechanism: no $X$, no transfer, no asymmetry. Second, the floor of the
benchmark window, $m_X>m_p$, until now part of a scaling assumption, is
\emph{derived} here as a survival condition. (The ceiling remains a benchmark.)
The running LHCb searches
(\texttt{Lb\_baryon\_lepton\_spectrum\_v1}, statistics in
$m(\Lambda_c^+\mu^-)$ and $m(p\mu^-)$;
\texttt{\_v2}, reach to $6.05\GeV$ through $m(\Xi_c^0\mu^-)$ and
$m(\Omega_c^0\mu^-)$) therefore probe the transfer sector of this paper.
[Derived; searches running.]

\paragraph{Freeze-out.}
With the threshold at $M_X$, equilibrium transfer shuts off by Boltzmann
suppression when $T$ falls below $M_X$; the standard logarithm gives
\begin{equation}
T_f \;\sim\; \frac{M_X}{25} \;\approx\; 200\ \mathrm{MeV}
\qquad (M_X=5\GeV).
\label{eq:tf}
\end{equation}
The alternative bookkeeping --- rate equals Hubble for the contact operator,
$T_f=(M_S^4/M_{\rm Pl})^{1/3}=47$~MeV --- lands within a factor four; the
closure below is insensitive to the choice. [Derived, two routes, one order.]

% ======================================================================
\section{The closure}
\label{sec:closure}

\subsection{The chain}
In equilibrium under the bias, the plasma carries
$n_B=b_w\,\chi\,\mu_w$ with $\chi\simeq T^2/6$ the baryon-number
susceptibility per relevant degree of freedom and $b_w$ the order-one transfer
coefficient collecting group-theory and channel counting. At freeze-out the
asymmetry is locked:
\begin{equation}
\eta_B \;=\; \frac{n_B}{n_\gamma}\bigg|_{T_f}
\;\simeq\; \frac{b_w\,(T_f^2/6)\,\mu_w}{0.244\,T_f^3}
\;\simeq\; \frac{b_w\,\mu_w}{1.5\,T_f}\,\Big.^{-1}
\quad\Longrightarrow\quad
b_w\,\mu_w \;\simeq\; 1.5\,\eta_B\,T_f .
\label{eq:closure}
\end{equation}

\subsection{The verdict, in numbers}
With $\eta_B=6.1\times10^{-10}$ (Planck) and the two freeze-out routes:
\begin{center}\small
\begin{tabular}{lcccc}
\toprule
route & $T_f$ & $b_w\mu_w$ required & $\mu_w$ (meV branch) & $b_w$ required\\
\midrule
threshold at $M_X$ & $200$~MeV & $1.8\times10^{-1}\eV$ & $34\meV$ & $\mathbf{5.4}$\\
$\Gamma=H$ & $47$~MeV & $4.3\times10^{-2}\eV$ & $34\meV$ & $\mathbf{1.3}$\\
\midrule
either & --- & --- & $\mu_w$ (TeV branch): $1.4\times10^{-16}\eV$ &
$\sim10^{15}$\\
\bottomrule
\end{tabular}
\end{center}
\textbf{On the meV branch the accounting closes with nothing forced:} the
required transfer coefficient is $b_w\simeq1$--$5$, in the same order-one class
as the $c_i$ of Article~IV. \textbf{On the TeV branch the mechanism fails by
fifteen orders} and no adjustment of $T_f$, $\chi$, or channel counting ---
each worth a factor of a few --- can bridge it. [Computed; script
\texttt{d4\_closure.py}.]

\subsection{What the verdict means}
The mechanism is neither confirmed nor decorative: it is \emph{conditional on
one microscopic derivation} --- the normalisation scale of the winding
coupling. If the winding is carried by the condensate phase, whose own scale is
the meV sector, the branch that closes is the natural one and this paper has
quantitative content. If the winding normalisation is instead set at the wall,
the mechanism is dead and \S\ref{sec:die}(iii) applies. The question is sharp,
answerable within the effective theory of Articles~II--III, and is filed as
the gate of this paper. [Open, and declared.]

% ======================================================================
\section{The direction and the amplitude, derived}
\label{sec:dipole}

Earlier drafts posited ``a fixed cosmological direction'' and a benchmark
$A\sim10^{-3}$. Both are derivable. A condensate defines a rest frame; a
directional asymmetry seen from Earth arises from \emph{our motion} through
it. If the condensate is comoving with the cosmic rest frame --- the natural
reading of Article~III's cosmology --- then
\begin{equation}
\hat n_{\rm DDF} \;=\; \hat n_{\rm CMB}\ \big[(l,b)\simeq(264^\circ,48^\circ)\big],
\qquad
A_{\rm DDF} \;=\; \beta \;=\; \frac{v_{\rm CMB}}{c} \;=\; 1.23\times10^{-3}.
\label{eq:dipole}
\end{equation}
The direction is thereby fixed \emph{a priori} --- no sky scan, no trials
factor --- and the amplitude, previously a benchmark, becomes the dipole
velocity. Any test of Prediction~I uses~(\ref{eq:dipole}) frozen before the
data are seen. [Derived under the comoving assumption, which is declared.]

% ======================================================================
% ---- Section: residual operator and SME dictionary (Phase A) ----
% The full finalized text of this section is maintained in
% proofs/section_dictionary_finale.tex and is \input here verbatim.
% ======================================================================
% Article 8 (DDF) — Section : The residual operator and its SME dictionary
% VERSION FINALE — corrections appliquées :
%   (1) dimension-7 : /Lambda^3 (et non /Lambda^2)
%   (2) les deux productions LHCb citées (v1 = statistique, v2 = Xi/Omega)
%   (3) hypothèse du pont de saveur déclarée (paragraphe dédié)
%   (4) une seule version de la section, un seul \label
%   (5) modulation sidérale horloges/comagnétomètres rétablie (ligne 5)
%   (6) langage « bounds, inserted » mis à l'indicatif
%   (7) règle de verdict forte : fenêtre à DEUX mâchoires
% ======================================================================

\section{The residual operator and its SME dictionary (Phase A)}
\label{sec:dictionary}

If the winding bias survives recombination, its four-dimensional imprint is a
background vector coupled to the currents the transfer operator carries. We
write it model-independently as
\begin{equation}
\Delta\mathcal{L}_{\rm res}
=
V_\mu\Big[\,q_B\,j_B^\mu+\sum_f q_f\,\bar\psi_f\gamma^\mu\psi_f\Big],
\label{eq:residual}
\end{equation}
with $V_\mu$ fixed in the cosmic rest frame. In the Earth frame,
$V_0\simeq\mu_w$ and $\vec V\simeq\mu_w\,\beta\,\hat n_{\rm CMB}$ with
$\beta=v_{\rm CMB}/c=1.23\times10^{-3}$ (\S\ref{sec:dipole}). Equation
(\ref{eq:residual}) is the object the experiments judge. Its translation into
the Standard-Model Extension (SME) is the following dictionary; bounds are
taken from the January~2026 edition of the Kosteleck\'y--Russell data tables
(arXiv:0801.0287, updated), cited here rather than the RMP~2011 edition.

\begin{center}\small
\begin{tabular}{p{3.8cm}p{2.3cm}p{3.9cm}p{4.0cm}}
\toprule
Piece of $V_\mu\times$ current & SME object & Observable channel & Bound \\
\midrule
Universal vector, constant, $q_f=q_B$ & $a_\mu^{B}$ & removable in flat space;
\emph{countershaded} by gravity (WEP) &
$\bar a_{\rm eff}\sim10^{-10}\text{--}10^{-11}$~GeV
[Kosteleck\'y--Tasson 2011; MICROSCOPE 2019] \\
Species differences & $\Delta a_0$ & clock comparisons &
[KR jan.~2026, proton sector D10--D14] \\
Axial contamination & $b_\mu^{p,n,e}$ & comagnetometers He/Xe &
$\sim10^{-33}\text{--}10^{-34}$~GeV \\
Gradient $\nabla V_0$ & --- & composition-dependent fifth force &
[E\"ot-Wash; neutron interferometry] \\
Spatial part, cosmic-fixed (matter sector) & sidereal harmonics &
clock and comagnetometer modulation & [KR jan.~2026, sidereal sector] \\
Spatial part, cosmic-fixed (flavor sector) & two-quark $\Delta a$ &
neutral-meson oscillations and decay asymmetries &
K: $\Delta a\sim10^{-18}$~GeV (KTeV/KLOE);
$B^0$: $\sim10^{-14}$~GeV (BaBar);
$B_s$: $\Delta a_\perp<1.2\times10^{-12}$~GeV,
$\Delta a_T-0.396\,\Delta a_Z\in(-0.8;\,3.9)\times10^{-13}$~GeV (D0) \\
Flavor-non-universal ($X$-mediated, four-fermion, dimension~7) &
off-diagonal, $\varepsilon^2$-fed & $B_s$--$\bar B_s$ mixing at second order &
protected; D0 applies at $\varepsilon^2$ [quantified in D2/D4] \\
\bottomrule
\end{tabular}
\end{center}

\paragraph{Method anteriority (R11).}
The observable of Prediction~I --- a beauty-hadron decay asymmetry correlated
with the sidereal phase --- is not a new invention: it is the established
observable class of the SME matter sector, measured and bounded by
D0~[arXiv:1506.04123], BaBar, KTeV and FOCUS, with ongoing LHCb charm analyses
[arXiv:2507.05457]. We cite this class as method anteriority: the prediction
borrows its method, and the class's null results are its judges.

\paragraph{Proposition A1 (removability, amended).}
\textit{In flat space, if the residual couples universally to the conserved
baryon current, no non-gravitational local observable --- including the
$b\to s\mu^+\mu^-$ imprint of Prediction~I --- survives.}
Under $\psi\to e^{iB_\psi a\cdot x}\psi$ with universal $B_\psi$, the constant
$a_\mu^B$ is removed; only species differences (clocks) and gradients (fifth
force) remain. \textbf{Amendment:} in the presence of gravity the removal is
incomplete --- the universal $a_\mu$ re-enters through matter--gravity couplings
(countershading, Kosteleck\'y--Tasson, PRD~83, 016013 (2011)), and is bounded by
weak-equivalence tests at $\bar a_{\rm eff}\sim10^{-10}\text{--}10^{-11}$~GeV
(MICROSCOPE). The proposition therefore reads \emph{no non-gravitational
observable}; the gravitational one is judged at $10^{-10}$~GeV, weak compared to
the meson bounds and not verdict-changing, but it closes the ``constant part is
unobservable'' escape. \textit{[Derived, field-redefinition identity; amendment
from the literature.]}

\paragraph{Proposition A2 (flavor obstruction).}
\textit{A surviving decay imprint requires non-universal flavor charges.}
Since $b$ and $s$ carry the same $B=1/3$, the universal redefinition of A1
leaves $\bar s\gamma^\mu b$ invariant and removes the imprint with it. A
directional term $P_1(\hat p_{\mu\mu}\cdot\hat n_{\rm DDF})$ can therefore
exist only if $V_\mu$ couples with charges that differ between flavors ---
precisely the structure the $X$ mediators of Article~VII supply
($B=1$, $L=1$, coupled to $c\,\mu$ channels). \textit{[Derived.]}

\paragraph{The flavor bridge (declared assumption).}
The $X$ states of Article~VII are established at the $c\,\mu$ vertex
($X^0_{(0,1,1)}\to\Lambda_c^+\mu^-$); the imprint of Prediction~I lives in
$b\to s\mu^+\mu^-$. The narrowing theorem below therefore carries a declared
assumption: \emph{that the flavor couplings of $X$ extend from the $c\mu$
vertex to the $bs$ current.} This extension is not automatic; it is to be
specified --- or refuted --- by the operator construction of D2. Every
statement of the form ``tests both at once'' holds under this assumption and
not otherwise. \textit{[Declared assumption.]}

\paragraph{The narrowing theorem (verdict of the numbers).}
For a two-quark vector insertion $\varepsilon\,(\bar s\gamma^\mu b)\hat V_\mu$
to imprint a directional modulation of order $\beta$ on $b\to s\mu^+\mu^-$, the
coefficient must sit at $\varepsilon\sim10^{-5}\text{--}10^{-6}$~GeV; the D0
bound on precisely this object is $10^{-12}\text{--}10^{-13}$~GeV.
\textbf{The strong version, realized through a two-quark insertion, is excluded
by six to seven orders of magnitude.} Proposition~A2, however, identified the
only exit, and the computation confirms it as the sole one: a four-fermion
operator mediated by the $X$ states of Article~VII,
$V_\mu(\bar s\gamma^\mu b)(\bar\mu\mu)/\Lambda^3$ --- of dimension seven, hence
the cubic scale --- feeds $B_s$--$\bar B_s$ mixing only at second order
($\varepsilon^2$), and thus evades the D0 bound structurally while producing
the directional imprint in the decay. The clause ``protected,
flavor-non-universal structure'' of the pre-registered verdict rule is
therefore promoted from a cautious clause to a \textbf{narrowing theorem}:
under the flavor-bridge assumption above, Prediction~I can survive only through
the $X$ operator of Article~VII. The two articles are no longer neighbours; one
is the survival condition of the other, and the ongoing LHCb searches
(\texttt{Lb\_baryon\_lepton\_spectrum\_v1}, which carries the statistics in
$m(\Lambda_c^+\mu^-)$ and $m(p\mu^-)$, and
\texttt{Lb\_baryon\_lepton\_spectrum\_v2}, which extends the reach through
$m(\Xi_c^0\mu^-)$ and $m(\Omega_c^0\mu^-)$ to 6.05~GeV) test both at once.
\textit{[Derived, structural; protection factor quantified in D2/D4.]}

\paragraph{Verdict rules (pre-registered, amended; the two-jaw window).}
The strong version does not survive by smallness alone: an $X$ operator small
enough to evade every bound while imprinting nothing is not a surviving
prediction but an unfalsifiable one. The strong version therefore lives in a
\emph{window with two jaws}, in the manner of the $\lambda$ window of
Article~III (astrophysical ceiling, cosmological floor):
\begin{itemize}
\item \textbf{Weak version.} If, at the couplings required by $\eta_B$, the
window below is empty --- in particular if the $\varepsilon^2$ contribution to
$B_s$--$\bar B_s$ mixing exceeds the D0 bound --- the residual is excluded as a
present-day coupling: the bias is purely primordial and Prediction~I is
withdrawn as a residual imprint.
\item \textbf{Strong version.} If, at the couplings required by $\eta_B$, the
$X$-mediated structure satisfies \emph{both} jaws ---
(\emph{i}) the decay-interference ratio remains of order one, so that the
imprint reaches $A_{\rm DDF}\simeq\beta$, and
(\emph{ii}) the $\varepsilon^2$ mixing contribution stays below the D0 bound
--- then Prediction~I is retained with $\hat n_{\rm DDF}=\hat n_{\rm CMB}$ and
$A_{\rm DDF}=\beta=1.23\times10^{-3}$ (\S\ref{sec:dipole}), and the same $X$
operator that baryogenesis needs becomes testable in the LHCb searches of
Article~VII.
\end{itemize}
Whether the window is non-empty at the couplings $\eta_B$ requires is the
single gate computation of this article (D4); $\chi_w$, $\Gamma_w$ and $\mu_w$
are its inputs.

\paragraph{Correction of an earlier conflation.}
The figures $10^{-31}$--$10^{-33}$~GeV apply to the \emph{axial} coefficients
$b_\mu$ (comagnetometers). The DDF residual is naturally \emph{vector}; its
judges are the clock, WEP, fifth-force and meson-oscillation lines of the
dictionary, whose bounds are strong but of a different order.

\paragraph{How this section dies.}
(i) If the clock and meson bounds of the dictionary --- now inserted ---
exclude the coefficient required at $A_{\rm DDF}\sim10^{-3}$ for every
non-universal charge assignment, the strong version dies and the weak version
is adopted. (ii) If the residual is shown to be exactly universal,
Propositions A1--A2 remove the imprint outright, with the same consequence.
(iii) If $\mu_w$ itself vanishes (\S\,failure conditions), the dictionary has
no object to judge. \textbf{(iv) If the two-jaw window is empty at the
couplings $\eta_B$ requires --- either because the $\varepsilon^2$ mixing
contribution exceeds the D0 bound, or because the decay-interference ratio
cannot reach order one without doing so --- the strong version dies, and with
it the only surviving route of Prediction~I.} (v) If the flavor-bridge
assumption is refuted by the operator construction of D2, the narrowing
theorem loses its object: the $X$ operator no longer reaches
$b\to s\mu^+\mu^-$, and Prediction~I is withdrawn regardless of the window.


% ======================================================================
\section{The judges}
\label{sec:judges}

\begin{center}\small
\begin{tabular}{p{4.4cm}p{9.2cm}}
\toprule
judge & what it constrains \\
\midrule
$\eta_B=6.1\times10^{-10}$ (Planck) & the target of the closure
(\S\ref{sec:closure}) \\
nucleon lifetimes ($\tau_p>10^{34}$~yr) & forces the lock of \S\ref{sec:lock};
any thresholdless transfer is dead on arrival \\
neutral-meson oscillations (D0, BaBar, KTeV) & the flavor sector of any
present-day residual (\S\ref{sec:dictionary}) \\
WEP / countershading (MICROSCOPE) & the universal constant part of the
residual \\
clock comparisons, comagnetometers & species differences; axial contamination \\
baryon isocurvature (Planck, few \%) & any spatial variation of $\mu_w$ at
recombination scales \\
BBN & the mechanism must be complete before $T\sim1$~MeV --- satisfied,
$T_f\gtrsim47$~MeV \\
LHCb searches \texttt{v1}/\texttt{v2} (running) & the $X$ transfer sector;
under the flavor bridge, Prediction~I \\
\bottomrule
\end{tabular}
\end{center}

% ======================================================================
\section{Epistemic status}
\label{sec:status}

\begin{center}\small
\begin{tabular}{p{7.2cm}p{6.4cm}}
\toprule
element & status \\
\midrule
mechanism class (spontaneous baryogenesis) & inherited, cited
[Cohen--Kaplan] \\
winding charge on the interval & derived (topology) \\
gradient $\partial_y\Phi$ = well slope $F$ & derived (Article~II) \\
normalisation scale $\kappa_w$ & \textbf{open --- the gate} \\
nucleon lock, $m_X>m_p$ & derived \\
$X$ states as necessary mediators & derived (given the lock) \\
freeze-out $T_f\sim0.05$--$0.2\GeV$ & derived, two routes \\
closure $b_w\simeq1$--$5$ (meV branch) & computed \\
frozen geometry at $T_f$ & assumption, declared \\
$\hat n_{\rm DDF}=\hat n_{\rm CMB}$, $A=\beta$ & derived under the comoving
assumption \\
flavor bridge ($c\mu\to bs$) & declared assumption (D2 decides) \\
two-jaw window of the strong version & pre-registered; second jaw to be
quantified (D2) \\
\bottomrule
\end{tabular}
\end{center}

% ======================================================================
\section{How this paper dies}
\label{sec:die}

\emph{(i)} If the winding is shown not to couple to the transfer operator ---
$\mu_w=0$ --- there is no bias and no paper.
\emph{(ii)} If no $B=1,\ L=1$ state exists above $m_p$ (the LHCb searches and
their successors bounding the window ever tighter), the lock of
\S\ref{sec:lock} cannot be satisfied and the transfer is dead by the nucleon
lifetime.
\emph{(iii)} If the winding normalisation is derived to sit at the wall scale,
the closure fails by fifteen orders and the mechanism is withdrawn.
\emph{(iv)} If the two-jaw window of \S\ref{sec:dictionary} is empty at the
couplings $\eta_B$ requires, the strong version of Prediction~I dies; the
weak version (purely primordial bias) survives only as far as (i)--(iii)
allow.
\emph{(v)} If the flavor-bridge assumption is refuted by the explicit operator
construction, Prediction~I is withdrawn regardless of the window.
\emph{(vi)} If the parent framework dies at $8.2\,\mu$m (Article~I), this
paper goes with it.

% ======================================================================
\begin{thebibliography}{99}
\bibitem{CohenKaplan} A.~G. Cohen and D.~B. Kaplan,
\emph{Thermodynamic generation of the baryon asymmetry},
Phys.\ Lett.\ B \textbf{199}, 251 (1987);
\emph{Spontaneous baryogenesis}, Nucl.\ Phys.\ B \textbf{308}, 913 (1988).
\bibitem{CKN} A.~G. Cohen, D.~B. Kaplan and A.~E. Nelson,
\emph{Progress in electroweak baryogenesis},
Ann.\ Rev.\ Nucl.\ Part.\ Sci.\ \textbf{43}, 27 (1993).
\bibitem{DDF1} H.~Boufourou, \emph{The Dark Dimension Framework I: The Stage},
this series (2026).
\bibitem{DDF2} H.~Boufourou, \emph{II: The Well and the Tower}, this series
(2026).
\bibitem{DDF3} H.~Boufourou, \emph{III: The Medium}, this series (2026).
\bibitem{DDF4} H.~Boufourou, \emph{IV: One Scale and its Shadow}, this series
(2026).
\bibitem{DDF5} H.~Boufourou, \emph{V: From Type~I$'$ Strings --- a Candidate
that Reaches the Window}, this series (2026).
\bibitem{DDF6} H.~Boufourou, \emph{VI: What the Candidate Geometry Implies},
this series (2026).
\bibitem{DDF7} H.~Boufourou, \emph{VII: Seven Charge Classes, and Two that are
Empty}, this series (2026).
\bibitem{KR} V.~A. Kosteleck\'y and N.~Russell, \emph{Data Tables for Lorentz
and CPT Violation}, January 2026 edition, arXiv:0801.0287.
\bibitem{KT} V.~A. Kosteleck\'y and J.~D. Tasson, \emph{Matter-gravity
couplings and Lorentz violation}, Phys.\ Rev.\ D \textbf{83}, 016013 (2011).
\bibitem{MICROSCOPE} H.~Pihan-Le Bars \emph{et al.}, \emph{New test of Lorentz
invariance using the MICROSCOPE space mission}, Phys.\ Rev.\ Lett.\ (2019);
P.~Touboul \emph{et al.}, Phys.\ Rev.\ Lett.\ \textbf{129}, 121102 (2022).
\bibitem{D0} V.~M. Abazov \emph{et al.} (D0), \emph{Search for violation of
CPT and Lorentz invariance in $B_s^0$ meson oscillations},
arXiv:1506.04123.
\bibitem{LHCbCharm} \emph{New limits on CPT-symmetry violation in charm
mesons}, arXiv:2507.05457.
\bibitem{Planck} Planck Collaboration, \emph{Planck 2018 results. VI.
Cosmological parameters}, Astron.\ Astrophys.\ \textbf{641}, A6 (2020).
\end{thebibliography}

\vspace{6pt}
\noindent\small AI-assistance: general-purpose AI tools were used in the
preparation of this work, as in the rest of the series.

\end{document}
